\documentclass{beamer}

\usepackage[russian]{babel}
\usepackage{fontspec}

\setmainfont{Liberation Serif}
\setsansfont{Liberation Sans}
\setmonofont{Liberation Mono}

\usepackage{amsmath}
\usepackage{graphicx}
\usepackage{hyperref}
\usepackage{algorithm}
\usepackage{algorithmic}

\usetheme{Madrid}
\usecolortheme{seahorse}

\title{Алгоритм поиска пути A* (A-star)}
\author{Тупикова Екатерина Викторовна}
\date{21.03.2026}

\begin{document}

% Титульный слайд
\begin{frame}
    \titlepage
\end{frame}

% Содержание
\begin{frame}{Содержание}
    \tableofcontents
\end{frame}

\section{Что такое A*}
\begin{frame}{Что такое A*?}
    \begin{itemize}
        \item \textbf{A*} — это алгоритм поиска кратчайшего пути в графе.
        \item Почти как алгоритм Дейкстры.
        \item Использует эвристику для ускорения поиска.
        \item Широко применяется в робототехнике, видеоиграх.
    \end{itemize}
\end{frame}

\begin{frame} {Что такое A*?}
	\begin{figure}
		\centering
		\includegraphics[width=0.9\linewidth]{Figure_1.png}
	\end{figure}
\end{frame}
\begin{frame} {Что такое A*?}
	\begin{figure}
	\centering
	\includegraphics[width=0.9\linewidth]{astar_fails.png}
	\end{figure}
\end{frame}

\section{Математическая основа}
\begin{frame}{Основная формула}
    Для каждой рассматриваемой вершины $n$ вычисляется функция стоимости:
    \Large
    \[ f(n) = g(n) + h(n) \]
    \normalsize
    \begin{block}{Компоненты формулы:}
        \begin{itemize}
            \item \textbf{$g(n)$} — стоимость пути (сумма рёбер) от начальной точки до вершины $n$.
            \item \textbf{$h(n)$} — эвристическая стоимости кратчайшего пути от $n$ до цели.
            \item \textbf{$f(n)$} — общая стоимость пути от вершины $n$ до цели.
        \end{itemize}
    \end{block}
\end{frame}

\section{Эвристические функции}
\begin{frame}{Популярные эвристики}
    Выбор $h(n)$ зависит от типа перемещения:
    \begin{itemize}
        \item \textbf{Расстояние Манхэттена}.
        \item \textbf{Расстояние Чебышева}.
        \item \textbf{Евклидово расстояние}.
    \end{itemize}
    Эвристика должна быть допустимой, т.е. не завышать реальную стоимость.
\end{frame}

\section{Алгоритм работы}
\begin{frame}{Механика алгоритма}
    Алгоритм использует два "списка":
    \begin{enumerate}
        \item \textbf{Open List (Список непосещённых вершин)}: очередь с приоритетом по $f(n)$.
        \item \textbf{Closed List (Список посещённых вершин)}: массив, где посещённые отмечены.
    \end{enumerate}
    
    \begin{exampleblock}{Логика цикла}
        На каждой итерации берем вершину с \textbf{минимальным $f(n)$}, заносим в список непосещённых вершин соседей с новыми весами, заносим текущую вершину в посещённую.
    \end{exampleblock}
\end{frame}

\begin{frame}
	\centering
	\Huge \textbf{Пример работы}
\end{frame}

\section{Пример работы}
\begin{frame}{Пример работы}
	\begin{figure}
	\centering
	\includegraphics[width=0.7\linewidth]{GOD_SAVE_GEMINI/Step0.png}
	\end{figure}
\end{frame}
\begin{frame}{Пример работы}
	\begin{figure}
		\centering
		\includegraphics[width=0.7\linewidth]{GOD_SAVE_GEMINI/Step1.png}
	\end{figure}
\end{frame}
\begin{frame}{Пример работы}
	\begin{figure}
		\centering
		\includegraphics[width=0.7\linewidth]{GOD_SAVE_GEMINI/Step2.png}
	\end{figure}
\end{frame}
\begin{frame}{Пример работы}
	\begin{figure}
		\centering
		\includegraphics[width=0.7\linewidth]{GOD_SAVE_GEMINI/Step3.png}
	\end{figure}
\end{frame}
\begin{frame}{Пример работы}
	\begin{figure}
		\centering
		\includegraphics[width=0.7\linewidth]{GOD_SAVE_GEMINI/End_alg.png}
	\end{figure}
\end{frame}
\begin{frame}{Пример работы}
	\begin{figure}
		\centering
		\includegraphics[width=0.9\linewidth]{GOD_SAVE_GEMINI/GodSaveGemini.png}
	\end{figure}
\end{frame}

\section{Пример применения}
\begin{frame}{D*}
	\begin{figure}
		\centering
		\includegraphics[width=0.5\linewidth]{D*/step1.png}
	\end{figure}
	\footnote{Источник: Stentz, A. (1994). CMU-RI-TR-94-37.}
\end{frame}
\begin{frame}{D*}
	\begin{figure}
		\centering
		\includegraphics[width=0.5\linewidth]{D*/step2.png}
	\end{figure}
	\footnote{Источник: Stentz, A. (1994). CMU-RI-TR-94-37.}
\end{frame}
\begin{frame}{D*}
	\begin{figure}
		\centering
		\includegraphics[width=0.5\linewidth]{D*/step3.png}
	\end{figure}
	\footnote{Источник: Stentz, A. (1994). CMU-RI-TR-94-37.}
\end{frame}

\begin{frame}{Заключение}
    \centering
    \Huge Спасибо за внимание! \\
    \vspace{1cm}
\end{frame}

\end{document}
